\section{Latar Belakang}

Dunia kerja saat ini memiliki tuntutan yang semakin kompetitif. Kondisi ini
mendorong perguruan tinggi untuk menghasilkan lulusan yang tidak hanya
menguasai pemahaman teoritis, tetapi juga memiliki pengalaman praktis yang
relevan dengan bidang keilmuannya, salah satunya melalui kegiatan magang
industri. Menurut Kementerian Pendidikan, Kebudayaan, Riset, dan Teknologi,
praktik kerja merupakan kegiatan pembelajaran di luar kampus yang dilaksanakan
melalui kerja sama dengan mitra industri untuk memberikan pengalaman kerja
nyata kepada mahasiswa~\cite{Kemendikbud2023}. Program ini berfungsi sebagai
jembatan antara teori akademik dan praktik profesional di lapangan.

Kegiatan magang berkontribusi dalam mengembangkan kemampuan komunikasi, kerja
sama tim, serta adaptasi terhadap lingkungan profesional. Program ini dirancang
untuk meminimalkan kesenjangan antara teori akademik dan kebutuhan industri,
sehingga terbentuk pemahaman yang lebih komprehensif terhadap realitas
lapangan. Selain meningkatkan keterampilan teknis (\textit{hard skills}) sesuai
bidang studi, kegiatan ini juga membentuk karakter profesional dan kedisiplinan
dalam menyelesaikan tanggung jawab yang diberikan~\cite{Arisandi2022}.

Magang industri turut membuka peluang untuk membangun jaringan profesional yang
mendukung pengembangan karier, sekaligus meningkatkan kesiapan kerja di bidang
teknologi informasi. Keterlibatan dalam lingkungan kerja yang dinamis
memberikan gambaran nyata mengenai bagaimana solusi teknologi diciptakan untuk
menyelesaikan permasalahan bisnis yang kompleks~\cite{Arisandi2022}.

Sebagai mahasiswa Program Studi Teknik Informatika, pengalaman praktis tersebut
secara spesifik diwujudkan dalam bidang rekayasa perangkat lunak. Salah satu
aspek krusial dalam rekayasa sistem modern saat ini adalah pengembangan
antarmuka pengguna (\textit{user interface}/UI) yang responsif, intuitif, dan
terintegrasi dengan sistem \textit{back-end}, sebagai komponen penting dalam
aplikasi kesehatan digital skala enterprise. Kualitas antarmuka pengguna secara
langsung memengaruhi efisiensi operasional tenaga medis dalam mengakses dan
mengelola informasi kesehatan pasien. Pengalaman pengembangan aplikasi berbasis
Flutter selama masa studi diimplementasikan pada produk digital nyata yang
digunakan secara luas di industri.

Perusahaan yang menjadi lokasi pelaksanaan magang adalah PT Inti Utama
Solusindo (bagian dari Pharos Group), yang merupakan sebuah perusahaan farmasi
nasional dengan sejarah panjang dalam industri kesehatan Indonesia. Pharos
Group didirikan pada 30 September 1971 dengan visi menjadi rekanan pilihan
dalam penyediaan berbagai solusi kesehatan, serta misi untuk berinovasi demi
masa depan yang lebih sehat melalui kolaborasi yang bermakna\cite{aboutpharos}.

Saat ini, perusahaan berfokus pada pengembangan HealthyOne, sebuah
\textit{healthcare intelligence app} yang ditujukan bagi mitra klinik
kesehatan. HealthyOne berfungsi sebagai sistem manajemen klinik terintegrasi
yang mencakup pengelolaan stok obat, penjadwalan konsultasi dokter, layanan
marketplace B2B dan B2C, hingga sistem kasir (\textit{Point of Sales}/POS).
Keterlibatan dalam proyek berskala besar ini memberikan pengaruh langsung
terhadap pengelolaan data kesehatan yang menuntut ketelitian dan keandalan
sistem yang tinggi.

Dalam ekosistem kesehatan digital, pengembangan HealthyOne memerlukan fondasi
\textit{front-end} yang \textit{solid} untuk menjamin akurasi dan efisiensi
kerja tenaga medis. Antarmuka aplikasi tidak hanya dituntut menarik secara
visual, tetapi juga harus responsif dan mampu menangani integrasi API yang
kompleks agar aliran informasi berjalan lancar. Transisi teknologi platform
aplikasi web dari ekosistem Flutter ke React (Next.js) yang dilakukan oleh PT
Inti Utama Solusindo merupakan langkah strategis perusahaan untuk
mengoptimalkan performa serta \textit{maintainability} kode dalam jangka
panjang. Transisi ini juga memberikan kesempatan studi kasus nyata terhadap
keputusan pemilihan \textit{tech stack} pada aplikasi skala enterprise.

Secara teknis, lingkungan kerja profesional di PT Inti Utama Solusindo
memberikan pengalaman nyata dalam penerapan \textit{Software Development Life
    Cycle} (SDLC) pada proyek berskala \textit{enterprise} dengan kompleksitas
tinggi, yang secara langsung melengkapi teori yang telah diperoleh selama masa
perkuliahan. Program Diktisaintek Berdampak ini memberikan kesempatan untuk
mempraktikkan proses pembaruan fitur, penanganan \textit{bug}, serta
pengelolaan \textit{source code} pada repositori perusahaan. Koordinasi antar
divisi bersama \textit{Project Manager} (PM) dan tim \textit{Quality Assurance}
(QA) mendorong kedisiplinan dalam dokumentasi dan \textit{version control}
sesuai standar industri.